%% General definitions
\documentclass{article} %% Determines the general format.
\usepackage{a4wide} %% paper size: A4.
\usepackage[utf8]{inputenc} %% This file is written in UTF-8.
%% Some editors on Windows cannot save files in UTF-8.
%% If there is a problem with special characters not showing up
%% correctly, try switching "utf8" to "latin1" (ISO 8859-1).
\usepackage[T1]{fontenc} %% Format of the resulting PDF file.
\usepackage{fancyhdr} %% Package to create a header on each page.
\usepackage{lastpage} %% Used for "Page X of Y" in the header.
                      %% For this to work, you have to call pdflatex twice.
\usepackage{enumerate} %% Used to change the style of enumerations (see below).

\usepackage{amssymb} %% Definitions for math symbols.
\usepackage{amsmath} %% Definitions for math symbols.

\usepackage{tikz}  %% Pagacke to create graphics (graphs, automata, etc.)
\usetikzlibrary{shapes,arrows,calc,positioning}
\usetikzlibrary{graphs}
\usetikzlibrary{automata} %% Tikz library to draw automata
\usetikzlibrary{arrows}   %% Tikz library for nicer arrow heads
\usetikzlibrary{positioning}    %% Tikz library for more positioning options

\usepackage{hyperref} %% hyperlinks

\usepackage{chessfss} % contained in debian package texlive-games

\usepackage{listings}   % include & display python code

\usepackage[edges]{forest}  % for search trees

\usepackage{algorithm} % for algorithms
\usepackage{algpseudocode} % for algorithms of pseudocode

\tikzset{WhiteSquare/.style = {shape=rectangle, minimum width=28, minimum
    height=28, anchor=south west}}
\tikzset{BlackSquare/.style = {shape=rectangle, fill=lightgray, minimum
    width=28, minimum height=28, anchor=south west}}

%% Left side of header
\lhead{\course\\\semester\\Exercise \homeworkNumber}
%% Right side of header
\rhead{\authorname\\Page \thepage\ of \pageref{LastPage}}
%% Height of header
\usepackage[headheight=36pt]{geometry}
%% Page style that uses the header
\pagestyle{fancy}

%% Triangle node definitions (tikzpicture)
\tikzset{
    triangleup/.style={
        regular polygon,
        regular polygon sides=3,
        draw,
        shape border rotate=0,
        align=center,
        minimum size=8mm
    },
    triangledown/.style={
        regular polygon,
        regular polygon sides=3,
        draw,
        shape border rotate=180,
        align=center,
        minimum size=8mm
    }
}

\newcommand{\authorname}{Benjamin Regez\\Luis Tritschler}
\newcommand{\semester}{Spring Semester 2026}
\newcommand{\course}{Foundations of Artificial Intelligence}
\newcommand{\homeworkNumber}{12}

\renewcommand{\thesection}{Exercise \homeworkNumber.\arabic{section}}


\begin{document}

\section{}
\begin{enumerate}
    \item iteration:
    \begin{center}
        \begin{tikzpicture}[
        every node/.style = {
            rectangle, rounded corners,
            draw, align=center,
            top color=white, bottom color=white!20
        }]
        \node[triangleup, align=center] (root) {\scriptsize 1\\14};
        \node[triangledown, align=center, below right=0.5cm and 2cm of root] (c2) {};
        \draw[->, dashed] (root) to (c2);
        \node[triangleup, align=center, below right=0.5cm and 1.5cm of c2] (c22) {};
        \draw[->, dashed] (c2) to (c22);
        \node[triangledown, align=center, below right=0.5cm and 1cm of c22] (c222) {};
        \draw[->, dashed] (c22) to (c222);
        \node[draw=blue, circle, accepting, align=center, below right=0.5cm and 0.5cm of c222] (c2222) {\color{blue}14};
        \draw[->, dashed] (c222) to (c2222);
        \end{tikzpicture}
    \end{center}
    \item iteration:
    \begin{center}
        \begin{tikzpicture}[
        every node/.style = {
            rectangle, rounded corners,
            draw, align=center,
            top color=white, bottom color=white!20
        }]
        \node[triangleup, align=center] (root) {\scriptsize 2\\15};
        % right subtree
        \node[triangledown, align=center, below right=0.5cm and 2cm of root] (c2) {};
        \draw[->, dashed] (root) to (c2);
        \node[triangleup, align=center, below right=0.5cm and 1.5cm of c2] (c22) {};
        \draw[->, dashed] (c2) to (c22);
        \node[triangledown, align=center, below right=0.5cm and 1cm of c22] (c222) {};
        \draw[->, dashed] (c22) to (c222);
        \node[circle, accepting, align=center, below right=0.5cm and 0.5cm of c222] (c2222) {14};
        \draw[->, dashed] (c222) to (c2222);
        % left subtree
        \node[triangledown, align=center, below left=0.5cm and 2cm of root] (c1) {16\\ \scriptsize 1};
        \draw[->] (root) to (c1);
        \node[triangleup, align=center, below right=0.5cm and 1.5cm of c1] (c12) {};
        \draw[->, dashed] (c1) to (c12);
        \node[triangledown, align=center, below right=0.5cm and 1cm of c12] (c122) {};
        \draw[->, dashed] (c12) to (c122);
        \node[circle, draw=blue, accepting, align=center, below right=0.5cm and 0.5cm of c122] (c1222) {\color{blue}16};
        \draw[->, dashed] (c122) to (c1222);
        \end{tikzpicture}
    \end{center}
    \newpage
    \item iteration:
    \begin{center}
        \begin{tikzpicture}[
        every node/.style = {
            rectangle, rounded corners,
            draw, align=center,
            top color=white, bottom color=white!20
        }]
        \node[triangleup, align=center] (root) {\scriptsize 3\\$14.\overline{6}$};
        % right subtree
        \node[triangledown, align=center, below right=0.5cm and 2cm of root] (c2) {14\\ \scriptsize 1};
        \draw[->] (root) to (c2);
        \node[triangleup, align=center, below right=0.5cm and 1.5cm of c2] (c22) {};
        \draw[->, dashed] (c2) to (c22);
        \node[triangledown, align=center, below right=0.5cm and 1cm of c22] (c222) {};
        \draw[->, dashed] (c22) to (c222);
        \node[circle, draw=blue, accepting, align=center, below right=0.5cm and 0.5cm of c222] (c2222) {\color{blue}14};
        \draw[->, dashed] (c222) to (c2222);
        % left subtree
        \node[triangledown, align=center, below left=0.5cm and 2cm of root] (c1) {16\\ \scriptsize 1};
        \draw[->] (root) to (c1);
        \node[triangleup, align=center, below right=0.5cm and 1.5cm of c1] (c12) {};
        \draw[->, dashed] (c1) to (c12);
        \node[triangledown, align=center, below right=0.5cm and 1cm of c12] (c122) {};
        \draw[->, dashed] (c12) to (c122);
        \node[circle, accepting, align=center, below right=0.5cm and 0.5cm of c122] (c1222) {16};
        \draw[->, dashed] (c122) to (c1222);
        \end{tikzpicture}
    \end{center}
    \item iteration:
    \begin{center}
        \begin{tikzpicture}[
        every node/.style = {
            rectangle, rounded corners,
            draw, align=center,
            top color=white, bottom color=white!20
        }]
        \node[triangleup, align=center] (root) {\scriptsize 4\\$12.5$};
        % right subtree
        \node[triangledown, align=center, below right=0.5cm and 2cm of root] (c2) {14\\ \scriptsize 1};
        \draw[->] (root) to (c2);
        \node[triangleup, align=center, below right=0.5cm and 1.5cm of c2] (c22) {};
        \draw[->, dashed] (c2) to (c22);
        \node[triangledown, align=center, below right=0.5cm and 1cm of c22] (c222) {};
        \draw[->, dashed] (c22) to (c222);
        \node[circle, accepting, align=center, below right=0.5cm and 0.5cm of c222] (c2222) {14};
        \draw[->, dashed] (c222) to (c2222);
        % left subtree
        \node[triangledown, align=center, below left=0.5cm and 2cm of root] (c1) {11\\ \scriptsize 2};
        \draw[->] (root) to (c1);
        \node[triangleup, align=center, below right=0.5cm and 1.5cm of c1] (c12) {};
        \draw[->, dashed] (c1) to (c12);
        \node[triangledown, align=center, below right=0.5cm and 1cm of c12] (c122) {};
        \draw[->, dashed] (c12) to (c122);
        \node[circle, accepting, align=center, below right=0.5cm and 0.5cm of c122] (c1222) {16};
        \draw[->, dashed] (c122) to (c1222);
        % left-left subtree
        \node[triangleup, align=center, below left=0.5cm and 1cm of c1] (c11) {\scriptsize 1\\6};
        \draw[->] (c1) to (c11);
        \node[triangledown, align=center, below right=0.5cm and 1cm of c11] (c112) {};
        \draw[->, dashed] (c11) to (c112);
        \node[circle, draw=blue, accepting, align=center, below right=0.5cm and 0.5cm of c112] (c1122) {\color{blue}6};
        \draw[->, dashed] (c112) to (c1122);
        \end{tikzpicture}
    \end{center}
\end{enumerate}

\newpage
\section{}
\begin{enumerate}[(a)]
    \item With $\epsilon$-greedy we can only get a chance of
    0.5 to reach the subtree of $n_3$, namely with $\epsilon
    = 1$. With softmax policy, none of the children will get
    the probability of 1. The same holds for UCB1.\\
    Furthermore, even if we could come up with a policy
    that guarantees that $n_3$ is chosen in the next iteration,
    $n_3$ still has nodes to be expanded, so the node marked
    in grey won't be chosen anyway.
    \item $\epsilon = 0.2 \ \implies \ 1 - \epsilon = 0.8$\\
    $\mathbb{P}(n_1) = 0.8$, 
    $\mathbb{P}(n_2) = 0.1$, 
    $\mathbb{P}(n_3) = 0.1$
    \item Probabilities to choose $n_1$, $n_2$ or $n_3$ using
    softmax tree policy with $\tau$ = 10:
    \begin{enumerate}[]
        \item $n_1$: $e^{8/10} \approx 2.226$
        \item $n_2$: $e^{3/10} \approx 1.350$
        \item $n_3$: $e^{6/10} \approx 1.822$
        \item $\implies \Sigma_{i=1}^3 n_i = 5.398$
        \item $\mathbb{P}(n_1) = \frac{2.226}{5.398} \approx \underline{0.412}$
        \item $\mathbb{P}(n_2) = \frac{1.350}{5.398} \approx \underline{0.250}$
        \item $\mathbb{P}(n_3) = \frac{1.822}{5.398} \approx \underline{0.338}$
    \end{enumerate}
    \item Probabilities to choose $n_1$, $n_2$ or $n_3$ using
    UCB1 tree policy:
    \begin{enumerate}[]
        \item $n_1$: B$(n_1) = \sqrt{\frac{2 \cdot \ln(9)}{6}} \approx 0.856$ 
        $\implies 8 + 0.856 = \color{blue}8.856\color{black}$
        \item $n_2$: B$(n_2) = \sqrt{\frac{2 \cdot \ln(9)}{1}} \approx 2.096$ 
        $\implies 3 + 2.096 = 5.096$
        \item $n_3$: B$(n_3) = \sqrt{\frac{2 \cdot \ln(9)}{2}} \approx 1.482$ 
        $\implies 6 + 1.482 = 7.482$
        \item $n_1$ maximizes, so it gets chosen as best move.
        \item $\implies \mathbb{P}(n_1) = 1, \ \mathbb{P}(n_2) = 0, \
        \mathbb{P}(n_3) = 0$
    \end{enumerate}
\end{enumerate}

\section*{Statement on scientific integrity}
We solved the entire sheet on our own.

\end{document}
